\section*{ВВЕДЕНИЕ}
\addcontentsline{toc}{section}{ВВЕДЕНИЕ}

Управление проектами как самостоятельная инженерная дисциплина оформилось во второй половине XX века. К 2000-м годам ее принципы вышли за пределы строительства и оборонных программ -- сегодня по проектным методикам работают команды разработки программного обеспечения, маркетинга, научных исследований и государственных программ. Численность проектных команд растет, география распределенная, а планы все чаще ведутся параллельно по нескольким направлениям одновременно.

С увеличением числа задач и участников держать план в почте и электронных таблицах становится невозможно. Сроки расходятся между файлами, ответственные забываются, изменения в зависимостях задач не доходят до тех, кого они касаются. Отсюда возникла отдельная категория программных средств -- информационные системы управления проектами. Они хранят данные о проектах, задачах, сроках и составе участников в одной базе, и команда работает с единой версией плана.

К началу 2024 года ситуация с такими системами в России заметно изменилась. Atlassian свернул деятельность в стране и принудительно отключил облачные сервисы Jira и Trello для российских аккаунтов. У серверных версий с февраля 2024 года прекращена поддержка. Asana ограничила бесплатный тариф двумя пользователями. Российские альтернативы развиваются, но ни одна из проанализированных в работе систем не сочетает в одной точке нативную диаграмму Ганта, поддержку всех четырех типов зависимостей сетевого планирования, русскоязычный интерфейс, бесплатный тариф для небольшой команды и устойчивый доступ из РФ. Системы, закрывающие функциональную часть, либо ушли с рынка, либо превратились в универсальные корпоративные платформы, где управление проектами оказалось среди десятков других модулей.

Небольшим командам, ведущим несколько параллельных программных проектов, такие платформы избыточны. Запуск занимает недели, обучение пользователей -- дни, и большую часть функций команда не использует. Простые же канбан-доски быстро упираются в потолок: как только в проекте появляются взаимосвязанные работы, сроки и разные уровни доступа, доски перестают покрывать задачу планирования. Эти соображения и определили направление работы.

Цель настоящей работы -- разработка веб-ориентированной программно-информационной системы для управления проектами, рассчитанной на небольшие команды и объединяющей в одном инструменте реестр проектов, диаграмму Ганта с поддержкой четырех типов зависимостей задач и разграничение прав участников по трем ролям.

Для достижения поставленной цели необходимо решить следующие задачи:
\begin{itemize}
\item провести анализ предметной области и сопоставить функциональные возможности существующих ИС~УП, доступных российским пользователям;
\item сформировать техническое задание на разработку, описать прецеденты использования системы;
\item спроектировать архитектуру системы, схему базы данных и программный интерфейс взаимодействия клиентской и серверной частей;
\item реализовать клиентскую и серверную части, диаграмму Ганта с интерактивным изменением дат и порядка задач, режим добавления связей и разграничение прав участников;
\item провести системное тестирование программной системы.
\end{itemize}

Структура и объем работы. Работа состоит из введения, четырех разделов основной части, заключения, списка использованных источников и 2 приложений. Текст выпускной квалификационной работы равен \formbytotal{lastpage}{страниц}{е}{ам}{ам}.

Во введении сформулирована цель работы, поставлены задачи разработки, описана структура работы и приведено краткое содержание каждого из разделов.

В первом разделе приводится анализ предметной области управления проектами: рассматриваются ключевые понятия, методологии управления проектами, типовая функциональность ИС~УП и сопоставляются восемь существующих решений, доступных российским пользователям. По результатам сравнения сформулированы требования к разрабатываемой системе.

Во втором разделе приведено техническое задание на разработку: основание для разработки, цель и назначение, функциональные требования, требования к интерфейсу, описаны прецеденты использования системы, требования к надежности и информационной безопасности, к программной и технической совместимости, к оформлению документации.

В третьем разделе описано техническое проектирование системы: общая характеристика проектного решения, выбор средств и технологий реализации, архитектурное построение системы, проект базы данных с приведением ER-модели и описанием структуры таблиц, описание программного интерфейса между клиентской и серверной частями.

В четвертом разделе рассмотрена реализация системы: спецификация программных компонентов клиентской и серверной частей и результаты системного тестирования.

В заключении приведены основные результаты работы, полученные в ходе разработки.

В приложении А представлен графический материал. В приложении~Б представлены фрагменты исходного кода.
